\input{preamble.tex}

\renewcommand{\thesection}{}
\renewcommand{\thesubsection}{\arabic{subsection}}
\renewcommand{\thesubsubsection}{\alph{subsubsection}}

\begin{document}

% Indicate your name and your student number here (e.g., A01xxx)
\header{Assignment 8}{Wira Azmoon Ahmad}{A0149286R}

%======================================
\section{Rent Division}

\subsection{}
Since $v_{i1}=v_{i2}=\dots=v_{in}$, and $\sum_{j=1}^n v_{ij}=r$, then
\[\forall j: v_{ij} = \frac{r}{n}\]
Any assignment $j=\sigma(i)$ is valued the same by player $i$, so
\[\forall \sigma : v_{i\sigma(i)}=\frac{r}{n}\]
Since the outcome is EF, for the $\sigma$ in such an outcome,
\[\forall j: v_{i\sigma(i)}-p_{\sigma(i)}\geq v_{ij}-p_j\]
\[\forall j: -p_{\sigma(i)}\geq-p_j\]
\[\forall j: p_{\sigma(i)} \leq p_j\]
The price paid by player $i$ in any envy-free outcome
will be the least amount among all players.

\newpage
\subsection{}
Suppose we have an EF outcome with players $i,i'$, rooms $j,j'$,
with assignment 
\[\sigma(i)=j, \sigma(i')=j'\]
Since the outcome is EF,
\[v_{i\sigma(i)}-p_{\sigma(i)}\geq v_{ij'}-p_{j'}\]
Since $v_{ij}=v_j$, and $\sigma(i)=j$,
\[v_{j}-p_{j}\geq v_{j'}-p_{j'}\]
Rearranging,
\begin{equation}
    p_j \leq v_j-v_{j'} +p_{j'}
\end{equation}
By symmetry,
\[p_{j'} \leq v_{j'}-v_{j} +p_{j}\]
Rearranging once more,
\begin{equation}
    p_j \geq v_j-v_{j'} +p_{j'}
\end{equation}
Combining $(1),(2)$, we have for any EF outcome with price vector $\vec{p}$,
\begin{equation}
    \forall j, j':p_j=v_j-v_{j'}+p_{j'}
\end{equation}
Clearly $v_j-v_{j'}=k$ is a constant for each pair of rooms $j,j'$.\\
Note also that $\sum_{j=1}^n p_j = r$.
From $(3)$, for any EF price vector $\vec{p}^*\neq \vec{p}$,
\[\forall j, j': p_j^*>p_j \Rightarrow p_{j'}^* > p_{j'}\]
\[\forall j, j': p_j^*<p_j \Rightarrow p_{j'}^* < p_{j'}\]
In either case where, $p_j^*\neq p_j$, we have 
\[\sum_{j=1}^n p_j^* \neq r\]
A contradiction. 
Thus $p_j^*=p_j$ for any EF price vectors $\vec{p},\vec{p}^*$,
and room $j$. The price vector for an EF outcome is unique.

\subsection{}
Suppose we have an EF outcome where $\sigma(i)=j$,
for some player $i$ and least valued room $j$.
By envy-freeness,
\[\forall k: v_{i\sigma(i)}-p_{\sigma(i)} \geq v_{ik} - p_k\]
\[\forall k: v_{ij}-p_{j} \geq v_{ik} - p_k\]
Rearranging,
\[\forall k: p_{j} \leq v_{ij}-v_{ik} + p_k\]
Since $v_{ij}\leq v_{ik}$,
\[v_{ij} - v_{ik} \leq 0\]
Then,
\[\forall k: p_{j} \leq v_{ij}-v_{ik} + p_k \leq p_k\]
The price of room $j$ is the lowest of all rooms in any EF outcome.

%======================================
\newpage
\section{Mechanism Design}

\subsection{}
\subsubsection{}
False.
If $k\leq n$, where $n$ is the number of players,
the payment for each player is $0$ since the welfare of other
players is the same whether or not each player is present.\\
Consider $k>n$. 
Suppose we have 2 successful bidders $j,j'$ for some optimal outcome $o^*$.
In the optimal outcome, the items are assigned to the top $k$ valuations
by player.
Then, if player $j$ is absent, there exists an optimal outcome $o_{-j}^*$
where every player that was assigned an item in $o^*$, including $j'$,
is present, except for $j$, who would be replaced by the player with
the top $(k+1)^{th}$ value, say player $j''$.
This then implies that
\begin{equation}
    v_{j'}(o^*)=v_{j'}(o_{-j}^*)
\end{equation}
Player $j$ would pay
\begin{equation}
\begin{aligned}
    \sum_{i\neq j}[v_i(o_{-j}^*) - v_i(o^*)]
    &= \sum_{i\notin \{j,j'\}} 
    [v_i(o_{-j}^*) - v_i(o^*)] + v_{j'}(o_{-j}^*) - v_{j'}(o^*)\\
    &= \sum_{i\notin \{j,j'\}} [v_i(o_{-j}^*) - v_i(o^*)]
    \quad\text{(from (4))}
\end{aligned}
\end{equation}
Similarly, there exists an optimal outcome $o_{-j'}^*$
where $j'$ is replaced by $j''$. Player $j'$ would pay
\begin{equation}
    \sum_{i\notin \{j,j'\}} [v_i(o_{-j'}^*) - v_i(o^*)]
\end{equation}
Since $o_{-j}^*,o_{-j'}^*$ assign the items to the same players
besides $j,j'$, we have that
\begin{equation}
    \sum_{i\notin \{j,j'\}} v_i(o_{-j}^*) 
    = \sum_{i\notin \{j,j'\}} v_i(o_{-j'}^*)
\end{equation}
Then, comparing amounts paid by players $j,j'$ from $(5),(6)$,
\[\sum_{i\notin \{j,j'\}} [v_i(o_{-j}^*) - v_i(o^*)]
= \sum_{i\notin \{j,j'\}} [v_i(o_{-j'}^*) - v_i(o^*)]\]
Players $j,j'$ both pay the same amount.
It must be that two successful bidders pay the same amount of money.

\subsubsection{}
True. Consider an auction with 2 players and 2 items , where
\[v_{11}=v_{22}=2, v_{12}=3, v_{21}=1\]
We have an optimal outcome $o^*$ where item 1 is assigned to player 1,
and item 2 is assigned to player 2.
$o_{-1}^*$ will still assign item 2 to player 2,
while $o_{-2}^*$ will assign item 2 to player 1.
Then, player 1 pays 0, and player 2 pays 
\[v_{12}-v_{11}=1>0\]
The players pay different prices.

\subsection{}
No. The joint utility of two bidders $j,j'$ would be
\begin{equation}
\begin{aligned}
    &\sum_i v_i(o')-\sum_{i\neq j} v_i(o_{-j}'')
    + \sum_i v_i(o')-\sum_{i\neq j'} v_i(o_{-j'}''')\\
    &=2\sum_i v_i(o')
    -\left(\sum_{i\neq j} v_i(o_{-j}'') 
    + \sum_{i\neq j'} v_i(o_{-j'}''')\right)
\end{aligned}
\end{equation}
where $o',o'',o'''$ are the respective outcomes, 
given the reports of bidders $j,j'$.
Consider then an auction with 3 players and 1 item, 
with player 1 and 2 cooperating, and the valuations
\[v_1=v_3=2,v_2=1\]
However, player 1 and 2 misreports their values, reporting
a value of 3 for the item instead.
In the optimal outcomes, $o^*,o_{-2}^*$ assigns the item to player 1 or 3,
and $o_{-1}^*$ to player 3.
The joint utility is
\[2v_1-(v_3+v_2)=4-(2+2)=0\]
With the misreported values however, 
we have $o',o_{-2}'''$ only assigning to player 1, 
and $o_{-1}''$ assigns to player 2.
The joint utility is 
\[2v_1-(v_1+v_2)=4-(1+2)=1\neq 0\]
The bidders stand to gain joint utilty by not being truthful.

\subsection{}
No. Consider 3 players and 1 item with valuations
\[v_1=1,v_2=2,v_3=3\]
With truthful reporting, player 2's utility is 0, since she does not win.
By reporting a valuation $v'_2>3$, she can win the bid for a utility of
\[v_2-v_1=1\]
since the third-highest bid will be $v_1=1$.
Clearly the auction is not truthful.

\end{document}